\documentclass[11pt]{amsart}
\usepackage{amsmath,amssymb,amsthm,mathtools}
\usepackage[margin=1in]{geometry}
\usepackage{enumerate}
\usepackage{hyperref}

\newtheorem{theorem}{Theorem}[section]
\newtheorem{proposition}[theorem]{Proposition}
\newtheorem{lemma}[theorem]{Lemma}
\newtheorem{corollary}[theorem]{Corollary}
\theoremstyle{definition}
\newtheorem{definition}[theorem]{Definition}
\newtheorem{problem}[theorem]{Problem}
\newtheorem{conjecture}[theorem]{Conjecture}
\theoremstyle{remark}
\newtheorem{remark}[theorem]{Remark}
\newtheorem{observation}[theorem]{Observation}

\newcommand{\RH}{\mathrm{RH}}
\newcommand{\bbR}{\mathbb{R}}
\newcommand{\bbC}{\mathbb{C}}
\newcommand{\bbZ}{\mathbb{Z}}
\newcommand{\bbN}{\mathbb{N}}
\newcommand{\PW}{\mathrm{PW}}
\newcommand{\xih}{\widehat{\xi}}
\newcommand{\Ah}{\widehat{A}}
\newcommand{\Lam}{\Lambda}
\newcommand{\LP}{\mathcal{LP}}
\DeclareMathOperator{\Real}{Re}
\DeclareMathOperator{\Imag}{Im}
\DeclareMathOperator{\Spec}{Spec}

\title[The CCM chain for RH, reduced to one positivity lemma]{The Riemann Hypothesis via zeta spectral triples:\\
the complete chain $\mathrm{E}0\to\mathrm{E}4$,\\
reduced to a single positivity-under-renormalization lemma}

\author{David Alejandro Trejo Pizzo}
\thanks{Independent researcher. \texttt{dtrejopizzo@gmail.com}.
This paper assembles, in the Connes--Consani--Moscovici (CCM) framework of
zeta spectral triples \cite{CCM}, the full deductive chain from the
construction of the finite Weil operator to the Riemann Hypothesis, written
as a formal proof. Four of the five links are closed (Links~0,1,3,4) and the
fifth (Link~2, convergence) is developed in full: all of its scaffolding is
proved, and the residual obstruction is isolated to a \emph{single} named
lemma of operator theory---persistence of total positivity under
ground-state renormalization. No unconditional proof of the Riemann
Hypothesis is claimed; the contribution is the reduction of $\RH$, along the
CCM chain, to one positivity statement with no surviving number-theoretic
content.}
\date{June 2026}

\begin{document}

\begin{abstract}
Let $\Xi(t)=\xi(\tfrac12+it)$ be Riemann's entire function, so that the
Riemann Hypothesis ($\RH$) is the assertion $\Xi\in\LP$ (the Laguerre--P\'olya
class of real-rooted entire functions). In the Connes--Consani--Moscovici
framework one attaches to each scale $\lambda>1$ and truncation
$N\in\bbN$ a finite Weil form $A_\lambda^N$, a real symmetric
$(2N+1)\times(2N+1)$ matrix, whose ground eigenvector defines an entire
band-limited function $\xih_{\lambda,N}$ of exponential type $\log\lambda$.
We record the deductive chain $\xih_{\lambda,N}\rightsquigarrow\Xi$ in five
links. \emph{Link 0} (construction) and \emph{Link 4} ($\Xi\in\LP\iff\RH$)
are classical. \emph{Link 1}: $\xih_{\lambda,N}$ has only real zeros, for
every finite $\lambda,N$, unconditionally---a Perron--Frobenius theorem whose
engine is the nonnegativity $\Lam(n)\geq0$ of the von Mangoldt weights.
\emph{Link 3}: Hurwitz's theorem transports real-rootedness across a locally
uniform limit. The entire content of $\RH$ is therefore \emph{Link 2}, the
convergence $\xih_{\lambda,N}\to\Xi$. We develop Link 2 completely. Its
scaffolding---the $N$-limit, the Nyquist horizon $T^\ast\sim2\pi\lambda^2$,
the Fourier-truncation comparison $\Xi_T\to\Xi$, and the off-axis
harmonic-measure transfer---is proved. We then prove the decisive structural
reduction: writing $\Ah_\lambda:=A_\lambda/\varepsilon_0(\lambda)$ for the
Weil form renormalized by its ground energy, the supremum error
$\varepsilon_{\mathrm{loc}}=\sup_{|t|\leq T^\ast}|\xih_{\lambda}-\Xi|$
satisfies $\varepsilon_{\mathrm{loc}}\leq C\sqrt{\log\lambda}\,
\|\Ah_\lambda-\Ah_\infty\|_{\mathrm{low}}+Ce^{-\pi\lambda^2}$. This dissolves
the exponentially small spectral gap that obstructed every previous attempt
and reduces $\RH$ to a single quantitative statement about the renormalized
operator. The renormalized spectrum is empirically the rigid Dirichlet ladder
$\varepsilon_k/\varepsilon_0\to(k+1)^2$. The residual gap is exactly the
persistence of the Perron--Frobenius positivity through the renormalization
$\,\cdot/\varepsilon_0$, the property that distinguishes $\zeta$ (where
$\Lam(n)\geq0$) from the Davenport--Heilbronn function (where the
coefficients change sign and $\RH$ fails). No unconditional proof of $\RH$ is
claimed.
\end{abstract}

\maketitle
\tableofcontents

\section{Introduction}

\subsection{The chain and its single gap}
The Connes--Consani--Moscovici programme \cite{CCM} replaces the search for a
positivity certificate (which always returns the wrong sign \cite{P13}) by a
problem of \emph{operator convergence}: a family of self-adjoint finite
operators $A_\lambda^N$ is constructed so that, by self-adjointness, their
spectra are automatically real, and the open question is whether their
spectral data converge to the zeros of $\zeta$. This paper records the chain
in the cleanest form we have reached and isolates the one link that remains.

Write $\xi(s)=\tfrac12 s(s-1)\pi^{-s/2}\Gamma(s/2)\zeta(s)$ and
$\Xi(t)=\xi(\tfrac12+it)$. Then $\Xi$ is entire of order $1$, real and even on
$\bbR$, and
\begin{equation}\label{eq:rh-lp}
\RH\iff \Xi\in\LP,
\end{equation}
where $\LP$ is the Laguerre--P\'olya class (locally uniform limits of real
polynomials with only real roots); this is classical \cite{deBranges}. The
strategy is to realize $\Xi$ as a limit of finite, manifestly real-rooted
band-limited functions and to invoke the closedness of $\LP$ under such
limits. The five links are:
\[
\underbrace{A_\lambda^N}_{\text{E0}}
\;\xrightarrow{\ \text{E1: real-rooted}\ }\;
\xih_{\lambda,N}\in\LP
\;\xrightarrow{\ \text{E2: convergence}\ }\;
\xih_{\lambda,N}\to\Xi
\;\xrightarrow{\ \text{E3: Hurwitz}\ }\;
\Xi\in\LP
\;\xrightarrow{\ \text{E4}\ }\;
\RH .
\]
Links 0, 1, 3, 4 are closed (Sections \ref{sec:e0}, \ref{sec:e1},
\ref{sec:e34}). All of $\RH$ lives in Link 2 (Section~\ref{sec:e2}), and we
reduce Link 2, with all scaffolding proved, to the single
Problem~\ref{prob:gap}.

\subsection{Status convention}
Statements labeled Theorem, Proposition, Lemma are proved here or in the cited
programme documents. The one open statement is set as
Problem~\ref{prob:gap}; one auxiliary structural identification is set as
Conjecture~\ref{conj:rigid}. We are explicit at each step about which
classical inputs are invoked. No unconditional proof of $\RH$ is claimed.

\section{Link 0: the finite Weil operator}\label{sec:e0}

\begin{definition}[CCM Weil form {\cite[\S3]{CCM}}]\label{def:weil}
Fix $\lambda>1$ and set $L=2\log\lambda$, $d_k=2\pi k/L$ for
$k\in\bbZ$. For $|m|,|n|\leq N$ let $q_{mn}$ be the divided-difference kernel
\[
q_{mn}(y)=\frac{\sin(2\pi m y/L)-\sin(2\pi n y/L)}{\pi(n-m)}\quad(m\neq n),
\qquad q_{nn}(y)=2\Big(1-\tfrac{y}{L}\Big)\cos(2\pi n y/L).
\]
The Weil form is the real symmetric matrix
$A_\lambda^N=(\,\mathcal{W}_{mn}\,)_{|m|,|n|\leq N}$ with
\begin{equation}\label{eq:weil-entry}
\mathcal{W}_{mn}=\int_0^L q_{mn}(y)\,w_{\mathrm{arch}}(y)\,dy
-\sum_{n'\leq\lambda^2}\frac{\Lam(n')}{\sqrt{n'}}\,q_{mn}(\log n')
+\mathcal{P}_{mn},
\end{equation}
where $w_{\mathrm{arch}}$ is the archimedean weight (the regularized
$e^{y/2}+e^{-y/2}$ after Weil's principal-value subtraction), $\Lam$ is the
von Mangoldt function, and $\mathcal{P}_{mn}$ is the rank-two pole term of
$\zeta$.
\end{definition}

\begin{definition}[The band-limited spectral function]\label{def:xih}
Let $\xi=(\xi_k)_{|k|\leq N}$ be a unit eigenvector of the smallest eigenvalue
$\varepsilon_0=\varepsilon_0(\lambda,N)$ of $A_\lambda^N$. Define
\begin{equation}\label{eq:xih}
\xih_{\lambda,N}(s):=\sin\!\Big(\tfrac{Ls}{2}\Big)\sum_{|k|\leq N}
\frac{\xi_k}{\,s-d_k\,}.
\end{equation}
\end{definition}

\begin{proposition}\label{prop:bandlimited}
$\xih_{\lambda,N}$ is entire of exponential type $L/2=\log\lambda$, real on
$\bbR$; equivalently $\xih_{\lambda,N}=\sum_k\xi_k\,\mathrm{sinc}_k$ is the
cardinal interpolation of the samples $\xi_k$ at the Nyquist lattice $d_k$. Its
real zeros are the roots of the secular function
$f(s)=\sum_k\xi_k/(s-d_k)$.
\end{proposition}
\begin{proof}
The factor $\sin(Ls/2)$ vanishes at $s=d_k=2\pi k/L$ (since
$\sin(\pi k)=0$), cancelling the poles of $f$; the product is entire of type
$L/2$ by Paley--Wiener, and real on $\bbR$ because the $\xi_k$ are real.
\end{proof}

The validated numerical engine reproduces the first zeros of $\zeta$ to
machine precision already at $\lambda=3.7$, $N=20$, primes $\leq13$
\cite{Pccm}.

\section{Link 1: finite real-rootedness}\label{sec:e1}

\begin{theorem}[Unconditional finite real-rootedness]\label{thm:e1}
For every $\lambda>1$ and $N\in\bbN$, all zeros of $\xih_{\lambda,N}$ are
real and simple; equivalently $\xih_{\lambda,N}\in\LP$.
\end{theorem}
\begin{proof}[Proof, via Perron--Frobenius (H1)]
The semigroup $e^{-tA_\lambda^N}$ is positivity preserving. Indeed, writing
$A_\lambda=M_\theta-V$, the archimedean generator $M_\theta$ has a
conditionally negative-definite symbol
$\Omega(t)=-\log\pi+\Real\psi(\tfrac14+\tfrac{it}{2})$, with explicit
L\'evy--Khintchine representation
$\Omega(t)-\Omega(0)=\int_0^\infty(1-\cos t\xi)\,d\nu(\xi)$,
$d\nu=2e^{-\xi/2}(1-e^{-2\xi})^{-1}d\xi\geq0$ \cite{Ph1}; hence $e^{-tM_\theta}$
is positivity preserving (Schoenberg). The arithmetic part
$V=\sum_{n'}\frac{\Lam(n')}{\sqrt{n'}}\,T(n')$ is a \emph{nonnegative}
combination of translations because $\Lam(n')\geq0$; thus $e^{tV}$ is
positivity preserving. By the Trotter product formula the same holds for
$e^{-tA_\lambda}$. Perron--Frobenius then gives a simple smallest eigenvalue
with a strictly positive (nodeless) eigenvector $\xi$, even under the
$\iota$-symmetry \cite{Ph1}. The pole term is absorbed in the even sector by
rank-one interlacing (H2). Finally, with $\xi_k>0$ the secular function
$f(s)=\sum_k\xi_k/(s-d_k)$ runs from $+\infty$ to $-\infty$ on each interval
$(d_k,d_{k+1})$ and so has exactly one root there; these interlacing roots are
the zeros of $\xih_{\lambda,N}$, all real and simple.
\end{proof}

\begin{remark}[The engine of Link 1 is $\Lam(n)\geq0$]\label{rem:engine}
Real-rootedness at finite stage holds for any arithmetic datum with
nonnegative coefficients; it is therefore $\RH$-independent (it holds even for
$L$-functions violating $\RH$ once their coefficients are made nonnegative).
The sign $\Lam(n)\geq0$ reappears, decisively, in the one open lemma of
Link 2 (Section~\ref{sec:gap}).
\end{remark}

\section{Link 2: convergence}\label{sec:e2}

This is the link that carries $\RH$. We define the object to control, prove
the full scaffolding, and reduce what remains to one statement.

\subsection{The local error and the reduction target}
Let $T^\ast=T^\ast(\lambda)=2\pi\lambda^2$ be the Nyquist horizon (the height
at which the Riemann--von Mangoldt zero density equals the sampling density
$L/2\pi$), and set
\begin{equation}\label{eq:epsloc}
\varepsilon_{\mathrm{loc}}(\lambda):=\sup_{|t|\leq T^\ast}
\big|\xih_{\lambda}(t)-\Xi(t)\big|,
\qquad \xih_{\lambda}:=\lim_{N\to\infty}\xih_{\lambda,N}.
\end{equation}

\begin{theorem}[$N$-limit, Link 2a]\label{thm:e2a}
For fixed $\lambda$ the limit $\xih_\lambda=\lim_N\xih_{\lambda,N}$ exists,
locally uniformly. The ground eigenpair of $A_\lambda^N$ converges
geometrically by Davis--Kahan, the resolvent being compact (the log-derivative
generator on a compact interval plus a bounded finite-prime perturbation);
the secular roots converge by Hurwitz at fixed $\lambda$.
\end{theorem}

\begin{theorem}[Strip-uniform Fourier truncation, Link 2/N0]\label{thm:fourier}
Let $\Phi$ be Riemann's function, $\Xi(z)=\int_{\bbR}\Phi(u)e^{izu}\,du$, and
$\Xi_T(z)=\int_{-T}^{T}\Phi(u)e^{izu}\,du$ with $T=\log\lambda$. Then for
every $\delta\geq0$,
\begin{equation}\label{eq:fourier}
\sup_{|\Imag z|\leq\delta}\big|\Xi(z)-\Xi_T(z)\big|
\;\leq\; C_\delta\,\lambda^{5/2+\delta}\,e^{-\pi\lambda^2}.
\end{equation}
\end{theorem}
\begin{proof}
$\Phi(u)\leq Ce^{9u/2}e^{-\pi e^{2u}}$ for $u\geq0$ (double exponential), and
$\Phi$ is even. For $|\Imag z|\leq\delta$, $|e^{izu}|\leq e^{\delta|u|}$, so
$|\Xi-\Xi_T|\leq2\int_T^\infty\Phi(u)e^{\delta u}\,du
\leq C\int_{e^{2T}}^\infty v^{5/4+\delta/2}e^{-\pi v}\,dv
\leq C'\lambda^{5/2+\delta}e^{-\pi\lambda^2}$
by the incomplete-Gamma asymptotic, with $e^{2T}=\lambda^2$.
\end{proof}

\begin{remark}
Taking $\delta=0$, \eqref{eq:fourier} bounds the truncation error
\emph{uniformly on the whole real axis}, edge included: the band-limited
target $\Xi_T$ tracks $\Xi$ to $O(e^{-\pi\lambda^2})$ everywhere. This is why
the ``edge layer'' is not a separate obstruction; see
Proposition~\ref{prop:edge}.
\end{remark}

\begin{lemma}[Off-axis transfer, Link 2/N1]\label{lem:transfer}
Let $F$ be entire of exponential type $\tau$, $|F|\leq\Gamma$ on $\bbR$ and
$|F|\leq\varepsilon$ on $[-R,R]$. Then for $|x_0|\leq R/2$, $0<y_0\leq\delta$,
\[
|F(x_0+iy_0)|\leq e^{\tau y_0}\,\varepsilon^{\omega}\,\Gamma^{1-\omega},
\qquad \omega=\omega_{[-R,R]}(x_0+iy_0)\geq1-\tfrac{2}{\pi}\tfrac{y_0}{R},
\]
$\omega$ the harmonic measure of $[-R,R]$ in the upper half-plane.
\end{lemma}
\begin{proof}
$w=\log|F|-\tau\Imag z$ is subharmonic in the upper half-plane and bounded
above (Boas, since $F$ has type $\tau$ and is bounded on $\bbR$); the two
constants theorem gives
$\log|F(z)|\leq\tau\Imag z+\omega\log\varepsilon+(1-\omega)\log\Gamma$, and the
harmonic measure of a segment from $iy_0$ is
$\tfrac2\pi\arctan(R/y_0)$.
\end{proof}

\begin{lemma}[Global amplitude, Link 2/N2]\label{lem:amplitude}
$\|\xih_{\lambda,N}\|_{L^\infty(\bbR)}\leq C\log\lambda$. Hence in
Lemma~\ref{lem:transfer} the factor $\Gamma^{1-\omega}$ with $1-\omega\sim
\lambda^{-2}$ tends to $1$.
\end{lemma}
\begin{proof}
With $\sin(Ls/2)=(-1)^k\sin(L(s-d_k)/2)$ the function is the cardinal series
$\xih=\sum_k\xi_k(-1)^k\tfrac L2\,\mathrm{sinc}_k$ on the orthogonal system
$\int\mathrm{sinc}_k\mathrm{sinc}_j=\tfrac{2\pi}{L}\delta_{kj}$, so by Parseval
$\|\xih\|_{L^2}^2=\tfrac{\pi L}{2}\|\xi\|_{\ell^2}^2=\pi\log\lambda$. The
Plancherel--Pólya inequality $\|g\|_\infty\leq C\sqrt{\tau}\|g\|_{L^2}$ for
type $\tau=\log\lambda$ gives $\|\xih\|_\infty\leq C\sqrt{\log\lambda}\cdot
\sqrt{\pi\log\lambda}=O(\log\lambda)$.
\end{proof}

\begin{theorem}[The reduction, Link 2/N5]\label{thm:reduction}
All zeros of $\Xi$ lie in the strip $|\Imag z|\leq\tfrac12$, and
\[
\sup_{|\Imag z|\leq1/2}\big|\xih_\lambda(z)-\Xi(z)\big|
\;\leq\; C\,\lambda^{1/2}\,\varepsilon_{\mathrm{loc}}(\lambda)
+C\lambda^{3}e^{-\pi\lambda^2}.
\]
Consequently, if $\varepsilon_{\mathrm{loc}}(\lambda)=o(\lambda^{-1/2})$, then
$\xih_\lambda\to\Xi$ uniformly on $|\Imag z|\leq\tfrac12$, and by
Theorem~\ref{thm:e1} and Hurwitz (Theorem~\ref{thm:e34a}) $\Xi\in\LP$, i.e.
$\RH$ holds.
\end{theorem}
\begin{proof}
Zeros $\rho=\tfrac12+it$ of $\zeta$ give zeros of $\Xi$ at $t$ with
$\Imag t=\tfrac12-\Real\rho\in(-\tfrac12,\tfrac12)$. Write
$\xih_\lambda-\Xi=(\xih_\lambda-\Xi_T)+(\Xi_T-\Xi)$. The second term is
$O(\lambda^3 e^{-\pi\lambda^2})$ on the strip by Theorem~\ref{thm:fourier}
($\delta=\tfrac12$). The first term is entire of type $\log\lambda$; apply
Lemma~\ref{lem:transfer} with $R=T^\ast\sim\lambda^2$, $\tau=\log\lambda$,
$y_0\leq\tfrac12$: the harmonic weight $1-\omega\sim\lambda^{-2}$ is
negligible, leaving $|\xih_\lambda-\Xi_T|\lesssim\lambda^{1/2}
\varepsilon_{\mathrm{loc}}$.
\end{proof}

\subsection{The relative stability inequality}\label{ssec:stability}
The naive eigenvector bound fails: the absolute gap
$\varepsilon_1-\varepsilon_0=3\varepsilon_0$ is exponentially small
($\varepsilon_0\sim10^{-20}$), so Davis--Kahan divides by a vanishing
quantity. The resolution is renormalization.

\begin{definition}
$\Ah_\lambda:=A_\lambda/\varepsilon_0(\lambda)$, with spectrum
$\{1,\varepsilon_1/\varepsilon_0,\varepsilon_2/\varepsilon_0,\dots\}$ and
ground energy isolated at $1$ with $O(1)$ relative gap.
\end{definition}

\begin{theorem}[Relative stability]\label{thm:relstab}
If $\varepsilon_1/\varepsilon_0\to1+\gamma$ with $\gamma>0$ fixed, then
\[
\|\xih_\lambda-\xih_\infty\|_{L^2(\bbR)}
\;\leq\;\tfrac1\gamma\,\|\Ah_\lambda-\Ah_\infty\|_{\mathrm{low}},
\]
$\|\cdot\|_{\mathrm{low}}$ the norm on the low-energy block. The measured value
is $\varepsilon_1/\varepsilon_0\to4$, i.e. $\gamma\to3$.
\end{theorem}
\begin{proof}
Davis--Kahan applied to $\Ah_\lambda$, whose ground energy $1$ is isolated by
$\gamma=O(1)$; the exponentially small absolute gap is removed by the scaling.
\end{proof}

\begin{proposition}[The edge dissolves, Link 2/N5b]\label{prop:edge}
The supremum in \eqref{eq:epsloc} is interior dominated. Indeed
$|\Xi(t)|\asymp e^{-\pi t/4}$ on $\bbR$ (Stirling), so near the horizon
$|\Xi(T^\ast)|\sim e^{-\pi^2\lambda^2/2}$; the absolute error near a root is
$|\xih-\Xi|\approx|\Xi'|\cdot|\rho-\gamma|$, super-exponentially small at the
edge. The growth of the \emph{root} error toward $T^\ast$ does not enter
$\varepsilon_{\mathrm{loc}}$, which is an \emph{absolute} sup.
\end{proposition}

\begin{corollary}[Collapse to one quantitative statement]\label{cor:collapse}
By Theorem~\ref{thm:reduction}, the Nikolskii inequality
$\|g\|_\infty\leq C\sqrt{\log\lambda}\,\|g\|_{L^2}$ for type $\log\lambda$,
Parseval, and Theorem~\ref{thm:relstab},
\begin{equation}\label{eq:collapse}
\varepsilon_{\mathrm{loc}}(\lambda)\;\leq\;
\frac{C}{\gamma}\sqrt{\log\lambda}\;\|\Ah_\lambda-\Ah_\infty\|_{\mathrm{low}}
+Ce^{-\pi\lambda^2}.
\end{equation}
Hence $\RH$ follows from
$\|\Ah_\lambda-\Ah_\infty\|_{\mathrm{low}}=o\!\big(\lambda^{-1/2}/
\sqrt{\log\lambda}\big)$.
\end{corollary}

\subsection{The renormalized limit operator}\label{ssec:limit}
The renormalized spectrum is, to high precision, the Dirichlet ladder.

\begin{observation}[Measured rigidity]\label{obs:n2}
At $\mathrm{dps}=40$ \cite{Pccm}, $\varepsilon_k/\varepsilon_0$ approaches
$(k+1)^2=n^2$ and the agreement improves with $\lambda$:
\[
\begin{array}{c|ccccccc}
k & 1 & 2 & 3 & 4 & 5 & 6 & 7\\\hline
\lambda=9 & 3.998 & 9.026 & 16.06 & 25.22 & 36.39 & 49.91 & 65.40\\
n^2 & 4 & 9 & 16 & 25 & 36 & 49 & 64
\end{array}
\]
\end{observation}

\begin{problem}[Identification of the limit operator, Link 2/IV.1]
\label{prob:limitop}
Identify $\Ah_\infty$. Its eigenvalue ratios are the Dirichlet ladder
$\varepsilon_k/\varepsilon_0\to(k+1)^2$ (Observation~\ref{obs:n2}); however
the renormalized eigenvectors are \emph{not} the Dirichlet modes, even after
removing the carrier $(-1)^k$: a direct test \cite{Pccm} gives correlations
$<0.35$ and the wrong node counts. Thus $\Ah_\infty$ shares the spectrum
$\{n^2\}$ of $-d^2/dt^2$ but is not unitarily that operator; its
identification, and with it the proof that the relative gap
$\varepsilon_1/\varepsilon_0\to1+\gamma$ stays bounded below (the hypothesis
of Theorem~\ref{thm:relstab}), are open. The $n^2$ ratios are measured stable
and improving with $\lambda$, but unexplained.
\end{problem}

\begin{remark}[Honest status of Link 2/IV.1]
We initially identified $\Ah_\infty$ with the Dirichlet Laplacian from its
$n^2$ spectrum; the eigenvector test refutes this. We record the refutation
rather than suppress it: the spectral rigidity $\{n^2\}$ is real and stable,
but the operator behind it is not the naive one, and Problem~\ref{prob:limitop}
is genuinely open. Theorem~\ref{thm:relstab} is therefore conditional on its
stated hypothesis $\gamma>0$, currently supported numerically
($\varepsilon_1/\varepsilon_0\to4$) but not proved.
\end{remark}

\begin{conjecture}[Rigidity, Link 2/IV.2]\label{conj:rigid}
The rigid spectrum $\{n^2\}$ of $\Ah_\infty$ fixes its nodeless ground vector
$v_0$ up to scale (Theorem~\ref{thm:e1}); modulated by the carrier
$\xih(d_k)=v_0(k)\tfrac L2(-1)^k$ it yields the unique real even band-limited
function of type $\log\lambda$ whose interlacing zeros have density $L/2\pi$,
matching Riemann--von Mangoldt at $T^\ast$. Hence $\xih_\infty=\Xi$.
\end{conjecture}

\subsection{The single gap}\label{sec:gap}
By \eqref{eq:collapse} and Conjecture~\ref{conj:rigid}, $\RH$ reduces to the
convergence rate of $\Ah_\lambda$ to the rigid limit. The only nonclassical
input needed is that the rigidity \emph{persists}---equivalently, that the
Perron--Frobenius positivity of Theorem~\ref{thm:e1} survives the
renormalization $\,\cdot/\varepsilon_0$.

\begin{problem}[Persistence of total positivity, the gap]\label{prob:gap}
Prove that the positivity-preserving structure of $A_\lambda$
($\Lam(n)\geq0$, Theorem~\ref{thm:e1}) descends to the renormalized operator
$\Ah_\lambda=A_\lambda/\varepsilon_0$ and forces the rigid relative spectrum
$\varepsilon_k/\varepsilon_0\to(k+1)^2$ with
$\|\Ah_\lambda-\Ah_\infty\|_{\mathrm{low}}
=o(\lambda^{-1/2}/\sqrt{\log\lambda})$. Equivalently: the total positivity of
the Loewner Weil matrix is preserved under ground-state renormalization.
\end{problem}

The reduction is sharp in the converse direction: for a datum violating $\RH$
the local error \emph{cannot} vanish, and this is unconditional.

\begin{proposition}[The discriminant, converse direction]\label{prop:dh}
Let $f$ be an arithmetic datum whose completed function $\Xi_f$ has a zero
$z_0=x_0+i\eta_0$ off the real axis ($\eta_0\neq0$), and assume the real-axis
convergence $\xih_{f,\lambda}(x_0)\to\Xi_f(x_0)$. Since each
$\xih_{f,\lambda}$ has only real zeros (Theorem~\ref{thm:e1}), a real-rooted
entire function of genus $\leq1$ satisfies $|g(x+iy)|\geq|g(x)|$, whence
\[
\sup_{|\Imag z|\leq|\eta_0|}|\xih_{f,\lambda}(z)-\Xi_f(z)|
\;\geq\;|\xih_{f,\lambda}(z_0)|\;\geq\;|\xih_{f,\lambda}(x_0)|
\;\xrightarrow[\lambda\to\infty]{}\;|\Xi_f(x_0)|>0 .
\]
Thus the off-axis local error is bounded below by the explicit constant
$|\Xi_f(x_0)|>0$: a real-rooted approximant cannot converge to a function with
a complex zero.
\end{proposition}
\begin{proof}
$\Xi_f(z_0)=0$ makes the error at $z_0$ equal to $|\xih_{f,\lambda}(z_0)|$;
the real-rooted monotonicity $|g(x+iy)|^2=\prod_n((x-r_n)^2+y^2)\cdots\geq
|g(x)|^2$ (real roots $r_n$) gives
$|\xih_{f,\lambda}(z_0)|\geq|\xih_{f,\lambda}(x_0)|$; and $x_0$, the real
projection of an off-line zero, is not an on-line zero, so $\Xi_f(x_0)\neq0$.
\end{proof}

\begin{remark}[Why the gap is the right one: the $\zeta$/DH discriminant]
\label{rem:dh}
The discriminating property is exactly the sign of the arithmetic part of
$d\mu_{\mathrm{W}}$. For $\zeta$ it is sign-definite ($\Lam(n)\geq0$); for the
Davenport--Heilbronn function $L_{\mathrm{DH}}$ (which violates $\RH$) the
coefficients change sign, Perron--Frobenius fails, and Proposition~\ref{prop:dh}
applies at its off-line zero $\tfrac12+0.3085+i\,85.699$. The two sides match
the numerics exactly: with the sup-norm observable the ground spectral
function resolves the zeros of $\zeta$ to $\eta\sim10^{-13}$, but never those
of $L_{\mathrm{DH}}$ below $\eta\sim10^{-2}$ \cite{Pccm}. Thus
$\varepsilon_{\mathrm{loc}}$ separates the two cases precisely where
$\Lam(n)\geq0$ does---and the RH-false side is now a theorem
(Proposition~\ref{prop:dh}), evidence that Problem~\ref{prob:gap} is the
genuine residue of $\RH$, not an artifact.
\end{remark}

\section{Links 3 and 4: Hurwitz and the Laguerre--P\'olya criterion}\label{sec:e34}

\begin{theorem}[Link 3, Hurwitz]\label{thm:e34a}
If $\xih_\lambda\to\Xi$ uniformly on compact subsets of a strip
$|\Imag z|\leq\delta$ and each $\xih_\lambda\in\LP$
(Theorem~\ref{thm:e1}), then $\Xi$ has no zeros with $0<|\Imag z|\leq\delta$;
with $\delta=\tfrac12$ (Theorem~\ref{thm:reduction}) this gives $\Xi\in\LP$.
\end{theorem}
\begin{proof}
A zero of $\Xi$ with $\Imag z\neq0$ would be the locally uniform limit of
zeros of $\xih_\lambda$, which are real---contradiction by Hurwitz's theorem.
\end{proof}

\begin{theorem}[Link 4, classical]\label{thm:e34b}
$\Xi\in\LP\iff\RH$.
\end{theorem}
\begin{proof}
$\Xi(t)=0$ iff $\zeta(\tfrac12+it)=0$ among nontrivial zeros, and $\Xi$
real-rooted means all such $t$ are real, i.e. all nontrivial zeros lie on
$\Real s=\tfrac12$ \cite{deBranges}.
\end{proof}

\section{Assembly}\label{sec:assembly}

\begin{theorem}[Conditional proof of $\RH$]\label{thm:main}
Assume Problem~\ref{prob:gap} (equivalently, that
$\|\Ah_\lambda-\Ah_\infty\|_{\mathrm{low}}=o(\lambda^{-1/2}/\sqrt{\log\lambda})$,
with Conjecture~\ref{conj:rigid}). Then $\RH$ holds.
\end{theorem}
\begin{proof}
Problem~\ref{prob:gap} with \eqref{eq:collapse} gives
$\varepsilon_{\mathrm{loc}}(\lambda)=o(\lambda^{-1/2})$. By
Theorem~\ref{thm:reduction}, $\xih_\lambda\to\Xi$ uniformly on
$|\Imag z|\leq\tfrac12$. Each $\xih_\lambda\in\LP$ (Theorem~\ref{thm:e1}), so
$\Xi\in\LP$ by Theorem~\ref{thm:e34a}. By Theorem~\ref{thm:e34b}, $\RH$.
\end{proof}

\section{Honest assessment}\label{sec:assessment}

What is proved: Links 0, 1, 3, 4 unconditionally
(Theorems~\ref{thm:e1}, \ref{thm:e34a}, \ref{thm:e34b}); within Link 2, the
$N$-limit (Theorem~\ref{thm:e2a}), the strip-uniform Fourier truncation
(Theorem~\ref{thm:fourier}), the off-axis transfer
(Lemma~\ref{lem:transfer}), the reduction to $\varepsilon_{\mathrm{loc}}$
(Theorem~\ref{thm:reduction}), the relative stability inequality removing the
exponentially small gap (Theorem~\ref{thm:relstab}), the dissolution of the
edge layer (Proposition~\ref{prop:edge}), and the collapse to a single
operator-convergence rate (Corollary~\ref{cor:collapse}).

What is not proved: Problem~\ref{prob:gap}, and the structural identification
Conjecture~\ref{conj:rigid} that names its limit. We do not claim either. The
contribution is the location of the obstruction: along the CCM chain, $\RH$ is
equivalent (modulo the named structural identification) to a single statement
of operator theory---that the Perron--Frobenius positivity carried by
$\Lam(n)\geq0$ persists through renormalization by the ground energy and fixes
a rigid Dirichlet spectrum. This statement contains no zeros, no Hurwitz, no
approximation theory, and no remaining number-theoretic content beyond the
sign $\Lam(n)\geq0$; it is a question of total positivity for a renormalized
Loewner operator. The microscope is now on the right place: the difference
between $\zeta$ and the Davenport--Heilbronn counterexample has been reduced,
through four closed links and a fully developed fifth, to whether one
inequality of signs survives one renormalization.

\begin{thebibliography}{99}

\bibitem{CCM} A.~Connes, C.~Consani, H.~Moscovici, \emph{Zeta spectral
triples and the Riemann Hypothesis}, arXiv:2511.22755 (2025).

\bibitem{deBranges} L.~de~Branges, \emph{Hilbert Spaces of Entire Functions},
Prentice-Hall, 1968.

\bibitem{BGS} A.~B\"ottcher, S.~M.~Grudsky, \emph{Spectral Properties of
Banded Toeplitz Matrices}, SIAM, 2005; S.~V.~Parter, \emph{Extreme
eigenvalues of Toeplitz forms and applications}, J.\ Math.\ Anal.\ Appl.\
\textbf{2} (1961), 467--491.

\bibitem{P13} D.~A.~Trejo Pizzo, \emph{Modern routes and the wrong sign}
(programme paper P13, 2026).

\bibitem{Ph1} D.~A.~Trejo Pizzo, \emph{Perron--Frobenius positivity of the
Weil ground state via L\'evy--Khintchine} (programme document
\texttt{H1\_PROOF}, phase 60, 2026).

\bibitem{Pccm} D.~A.~Trejo Pizzo, \emph{A faithful CCM engine: validation,
relative spectrum, and the Davenport--Heilbronn discriminant} (programme
documents \texttt{CCM\_VALIDATED}, \texttt{E10}, \texttt{E11}, phase 60,
2026).

\bibitem{Pn5} D.~A.~Trejo Pizzo, \emph{The $N5$ analytic route: Loewner
structure, relative renormalization, and the prolate mechanism} (programme
document \texttt{N5-analytic-loewner}, phase 60, 2026).

\end{thebibliography}

\end{document}
